\documentclass[10pt,a4paper]{amsart}
\usepackage[margin=19mm]{geometry}
\usepackage{amsmath,amssymb,mathtools}
\usepackage[scr=rsfs,cal=euler]{mathalfa}
\usepackage{xfrac}
\usepackage[unicode,hidelinks,bookmarks=false]{hyperref}
\newcommand{\ie}{i.e.\ }
\newcommand{\cf}{cf.\ }
\newcommand{\resp}{respectively\ }
\newcommand{\Subsection}{\subsection}
\providecommand{\nobreakdash}{\nobreak\hbox{-}}
\setlength{\emergencystretch}{2em}
\multlinegap0pt
\usepackage{amssymb}
\usepackage{enumerate}
\usepackage{tikz}
\newtheorem{lemma}{Lemma}
\newcommand{\C}{\mathbb C}
\def\b1{\mathbf{1}}
\def\ba{\mathbf{a}}
\def\bo{\mathbf{0}}
\def\bs{\mathbf{s}}
  \def\delta{delta}%
  \def\gl{gl}%
\def\mm{\mathfrak{m}}
\title{Classification of Simple Cuspidal Modules over Nongraded Witt Lie Algebras}
\author{Genqiang Liu, Xiaoyao Zheng, Yufang Zhao}
\date{Source: arXiv:2607.28114v1 (CC BY 4.0)}
\begin{document}
\maketitle

\noindent\emph{Source and licence.} Classification of Simple Cuspidal Modules over Nongraded Witt Lie Algebras. Version arXiv:2607.28114v1 (CC BY 4.0), distributed by arXiv under the Creative Commons Attribution 4.0 International licence, http://creativecommons.org/licenses/by/4.0/, as recorded in the arXiv licence metadata for this exact version and shown on the abstract page https://arxiv.org/abs/2607.28114v1. Full licence notice: \texttt{licenses/arxiv-2607.28114-CC-BY-4.0.txt}.\par\smallskip

\noindent\textit{Editorial scope.} Counted unit: the article's lemma lem:gln-quotient (source0.tex lines 1471--1486). The article's complete original proof of it is delivered with it (source0.tex lines 1488--1683, one \texttt{proof} environment of 196 lines). No mathematics is written, repaired or re-derived: the statement and the proof are the article's own lines, in the article's own notation, and no retained line is substituted or re-flowed.  Retained uncounted as the source-local context the proof needs: lines 1063--1071; lines 1077--1083; lines 1348--1352; lines 1356--1360.  Nothing was omitted from any retained span, so the package carries no omission record.  No retained line contains a citation, so the wrapper carries no bibliography and no bibliography entry.  No retained line contains a figure environment, an image-inclusion command or drawing code, so no image is delivered and the package declares no asset dependency.  Editorial formatting only, in addition: an AMS \texttt{amsart} preamble that restates the article's own declarations and macros used by the retained lines, namely the environments \texttt{lemma} on separate counters, together with the standard packages those lines need.  The retained environments print in this document as Lemma 1 in order of appearance, whereas the article prints the same items with the numbering of its own full document.  The complete native source of the article as distributed in the arXiv e-print for this exact version is bundled unchanged as \texttt{sources/206386/source0.tex}.
\begin{lemma}\label{Hu}
Let \(V\) be finite-dimensional, and let
\(L\subseteq\gl(V)\) be a nonzero Lie algebra acting irreducibly on
\(V\). Then
\[
L=[L,L]\oplus Z(L),
\]
where \([L,L]\) is semisimple and \(\dim Z(L)\leq1\).
\end{lemma}

\begin{lemma}\label{lem:finite-jet}
If \(V\) is a finite-dimensional \(\mathfrak m_{\b1,\bo}\Delta\)-module, then
there exists \(N\geq 2\) such that
\[
\mm_{\b1,\bo}^N\Delta V=0.
\]
\end{lemma}

\begin{equation}\label{eq:basic-vector-field-bracket}
[fd_i,gd_j]
 =
 f\,d_i(g)d_j-g\,d_j(f)d_i.
\end{equation}

\begin{equation}\label{eq:filtration-bracket}
[\mathfrak m^p\Delta,\mathfrak m^q\Delta]
\subseteq
\mathfrak m^{p+q-1}\Delta
\end{equation}

\begin{lemma}\label{lem:gln-quotient}
If \(V\) is a finite-dimensional simple
\(\mathfrak m_{1,0}\Delta\)-module, then
\[
\mathfrak m_{1,0}^{2}\Delta V=0
\]
and
\[
(t_k-1-x_k)d_lV=0
\qquad
\text{for all }1\leq k,l\leq n.
\]
Consequently, the action of \(\mathfrak m_{1,0}\Delta\) on \(V\)
factors through \(\mathfrak{gl}_n\), and \(V\) is a simple
\(\mathfrak{gl}_n\)-module.
\end{lemma}

\begin{proof}
Let
\(
\rho\colon L\rightarrow\mathfrak{gl}(V)
\)
be the representation associated with \(V\).

By Lemma~\ref{lem:finite-jet}, there exists an integer \(N\geq2\)
such that
\begin{equation}\label{eq:finite-jet-property}
\mathfrak m^N\Delta V=0.
\end{equation}
Equivalently,  \(V\) is a module over the  quotient $
L/\mathfrak m^N\Delta$.


We first prove that
$
JV=\mathfrak m^2\Delta V=0$.
Let \(J^{(1)}=J\) and define inductively $
J^{(q+1)}=[J,J^{(q)}]$.
It follows from \eqref{eq:filtration-bracket}, by induction on \(q\),
that
\[
J^{(q)}\subseteq\mathfrak m^{q+1}\Delta.
\]
Together with \eqref{eq:finite-jet-property}, this implies that
\(\rho(J)\) is a nilpotent, and hence solvable, Lie algebra.

Since \(J\) is an ideal of \(L\), its image \(\rho(J)\) is an ideal
of $
G:=\rho(L)$.
The \(G\)-module \(V\) is irreducible. Therefore, by Lemma \ref{Hu},
\[
G=[G,G]\oplus Z(G),
\]
where \([G,G]\) is semisimple, and $Z(G)$ is equal to
the solvable radical of $G$. So as a  solvable ideal of $G$, we have
$\rho(J)\subseteq Z(G)$.
In particular,
\begin{equation}\label{eq:LJ-zero}
\rho([L,J])=0.
\end{equation}



We first eliminate the pure \(x\)-terms. Set
$ E_x=\sum_{i=1}^n x_id_i\in L$.
If $f\in\C[x_1,\ldots,x_n]$ is homogeneous of degree \(r\geq2\),
then
\[
[E_x,fd_j]=(r-1)fd_j.
\]
Since \(fd_j\in J\), equation~\eqref{eq:LJ-zero} gives
\(\rho(fd_j)=0\). Hence \(\mathfrak m_x^2\Delta V=0\), where
 $\mm_x
 =
 \langle x_1,\ldots,x_n\rangle
 \subseteq
 \mathbb C[x_1,\ldots,x_n]$.

For the pure \(t\)-terms, let \(y_i=t_i-1\),
\(\partial_i=t_i^{-1}d_i\), and
$E_t=\sum_{i=1}^n y_i\partial_i\in L$.
If \(f\in\C[y_1,\ldots,y_n]\) is homogeneous of degree \(r\geq2\),
then
\[
[E_t,f\partial_j]=(r-1)f\partial_j.
\]
Thus every such \(f\partial_j\) acts trivially on $V$. Let $
\mathfrak m_t
 =
 \langle t_1-1,\ldots,t_n-1\rangle
 \subseteq
 \mathbb C[t_1^{\pm1},\ldots,t_n^{\pm1}]
$. Modulo
\(\mathfrak m_t^N\), each Laurent monomial in the \(t_i\) has a
finite Taylor expansion in the \(y_i\). Together with
\eqref{eq:LJ-zero}, this shows that every element of
\(\mathfrak m_t^2\Delta\) acts trivially. Therefore
\begin{equation}\label{eq:pure-terms-zero}
\mathfrak m_t^2\Delta V=0,
\qquad
\mathfrak m_x^2\Delta V=0.
\end{equation}

We next prove that all mixed terms act trivially. More precisely,
we show that
\begin{equation}\label{eq:mixed-zero}
y^\ba x^\bs d_lV=0
\end{equation}
whenever \(|\ba|\geq1\) and \(|\bs|\geq1\).

By \eqref{eq:finite-jet-property}, the assertion is true when
\(|\ba|+|\bs|\geq N\). We now argue by descending induction on the total
degree
$q=|\ba|+|\bs|$,
and, for a fixed $q$, by ascending induction on \(|\ba|\).

The case \(|\ba|=0\) follows from
\(\mathfrak m_x^2\Delta V=0\). Suppose that \(|\ba|\geq1\).
For any \(1\leq i\leq n\), formula
\eqref{eq:basic-vector-field-bracket} gives
\begin{align*}
[y^\ba d_i,x^{\bs+e_i}d_l]
={}&
(s_i+1)y^\ba x^\bs d_l\\
&-
a_l y^{\ba-e_l}x^{\bs+e_i}d_i
-
a_l y^\ba x^{\bs+e_i}d_i.
\end{align*}
The second element \(x^{\bs+e_i}d_l\) in the bracket
 belongs to \(\mathfrak m_0^2\Delta\), and therefore
acts trivially by \eqref{eq:pure-terms-zero}. Hence the left-hand
side acts trivially on \(V\).
The term
$y^{\ba-e_l}x^{\bs+e_i}d_i$
has the same total degree \(q\), but has \(y\)-degree \(|\ba|-1\);
it acts trivially by the induction hypothesis on \(|\ba|\).
The term $y^\ba x^{\bs+e_i}d_i$
has total degree \(q+1\), and acts trivially by descending induction
on \(q\). Therefore
$(s_i+1)y^ax^sd_lV=0$.
Since \(s_i+1\neq0\),  \eqref{eq:mixed-zero} follows.

Every coefficient in \(\mathfrak m^2\), modulo \(\mathfrak m^N\),
is a linear combination of monomials of the following three types:
\[
y^\ba,\quad |\ba|\geq2;
\qquad
x^\bs,\quad |\bs|\geq2;
\qquad
y^\ba x^\bs,\quad |\ba|,|\bs|\geq1.
\]
The first two types act trivially by
\eqref{eq:pure-terms-zero}, and the third type acts trivially by
\eqref{eq:mixed-zero}. Together with
\eqref{eq:finite-jet-property}, this proves
\begin{equation}\label{eq:J-zero}
JV=\mathfrak m^2\Delta V=0.
\end{equation}

Thus \(V\) is naturally a simple module over \(L/J\). Define
\[
I=
\operatorname{span}_{\mathbb C}
\left\{
(t_k-1-x_k)d_l+J
\ \middle|\
1\leq k,l\leq n
\right\}
\subseteq L/J.
\]
From $$\aligned \
&[(t_i-1)d_j+\mm^2\Delta,(t_k-1-x_k)d_l+\mm^2\Delta]\\
&=\delta_{jk}(t_i-1)(t_k-1)d_l
-\delta_{li}(t_k-1-x_k)t_id_j+\mm^2\Delta\\
&=-\delta_{li}(t_k-1-x_k)d_j+\mm^2\Delta
\endaligned $$
and
$$\aligned \
&[x_id_j+\mm^2\Delta,(t_k-1-x_k)d_l+\mm^2\Delta]\\
&=\delta_{jk}x_i(t_k-1)d_l
-\delta_{li}(t_k-1-x_k)d_j+\mm^2\Delta\\
&=-\delta_{li}(t_k-1-x_k)d_j+\mm^2\Delta,
\endaligned $$
it follows that $[L/J,I]=I$ and $[I,I]=0$.

Let
$\overline{\rho}\colon L/J\rightarrow\mathfrak{gl}(V)$
be the induced representation and
$\overline{G}=\overline{\rho}(L/J)$.
Since \(I\) is an abelian ideal, \(\overline{\rho}(I)\) is a solvable
ideal of the reductive Lie algebra \(\overline{G}\). Hence
\[
\overline{\rho}(I)\subseteq Z(\overline{G}).
\]
Then we obtain
$\overline{\rho}(I)
 =
 \overline{\rho}([L/J,I])
 =
 [\overline{G},\overline{\rho}(I)]
 =
 0$.
Consequently,
$(t_k-1-x_k)d_lV=0$ for all $1\leq k,l\leq n$.
Finally, the kernel of the homomorphism \(\pi_2\) is
$ I=
\operatorname{span}_{\mathbb C}
\{Y_{ki}-X_{ki}\mid 1\leq k,i\leq n\}$ .
Therefore, $\frac{L/J}{I}\cong\mathfrak{gl}_n$.
The action of \(L\) on \(V\) factors through this quotient, so \(V\)
is a \(\mathfrak{gl}_n\)-module.
\end{proof}

\end{document}
